\documentclass{amsart}

% 数学相关宏包
\usepackage{amsmath,mathtools,amsthm,amssymb}
\usepackage{bbm}
\usepackage{mathrsfs}
\usepackage{centernot}

% 图形和绘图
\usepackage{graphicx,float}
\usepackage{tikz,tikz-cd}
\usepackage{pgfplots}
\usepackage[framemethod=TikZ]{mdframed}
\usetikzlibrary{decorations.pathmorphing,positioning,arrows,automata,calc,shapes,patterns}
\pgfplotsset{compat=1.18}
\usepgfplotslibrary{statistics}

% 表格和列表
\usepackage{booktabs,multirow,colortbl}
\usepackage{enumitem}
\setlist[itemize,enumerate,description]{noitemsep}
\setlistdepth{9}

% 算法
\usepackage{algorithm,algorithmic}

% 字体和颜色
\usepackage{xcolor,listings}
\usepackage{xeCJK}
\setCJKmainfont{Noto Serif CJK SC}

% 页面和链接
\usepackage{fancyhdr,caption}
\usepackage{hyperref}
\usepackage{cleveref}

% 工具宏包
\usepackage{etoolbox,letltxmacro,docmute}
\usepackage{gensymb}

% 自定义设置
\renewcommand{\arraystretch}{1.5}
\let\originalmathbb\mathbb
\renewcommand{\mathbb}[1]{%
    \ifstrequal{#1}{1}{\mathbbm{#1}}{\originalmathbb{#1}}%
}

% 颜色定义
\definecolor{lightblue}{RGB}{173,216,230}
\definecolor{lightgreen}{RGB}{144,238,144}
\definecolor{lightyellow}{RGB}{255,255,224}
\definecolor{lightgray}{RGB}{211,211,211}

		% 超链接设置
		\hypersetup{
		    colorlinks=true,
		    linkcolor=red,
		    filecolor=magenta,
		    urlcolor=cyan,
		    pdftitle={\@title},
		    pdfauthor={\@author},
		    pdfsubject={数学笔记},
		    pdfkeywords={数学},
		    bookmarksnumbered=true,
		    bookmarksopen=true,
		    bookmarksopenlevel=6,
		    pdfstartview=FitH,
		    pdfpagemode=UseOutlines,
		    pdfborder={0 0 0}
		}

% 定理环境
\theoremstyle{plain}
\newtheorem{theorem}{Theorem}[section]
\newtheorem{lemma}[theorem]{Lemma}
\newtheorem{corollary}[theorem]{Corollary}
\newtheorem{proposition}[theorem]{Proposition}

\theoremstyle{definition}
\newtheorem{definition}[theorem]{Definition}
\newtheorem{example}[theorem]{Example}
\newtheorem{exercise}[theorem]{Exercise}
\newtheorem{problem}[theorem]{Problem}

\theoremstyle{remark}
\newtheorem{remark}[theorem]{Remark}
\newtheorem{note}[theorem]{Note}
\newtheorem{observation}[theorem]{Observation}

% 解答环境
\newtheorem*{solution}{Solution}
\newtheorem*{proof*}{Proof}

% 图片路径
\graphicspath{{F:/资料/2025-summer/figures/}}

\author{乐绎华, Yue Yihua}
\address{中山大学, Sun Yat-sen University}
\email{yueyh@mail2.sysu.edu.cn}
\date{\today}
\title{Mathematical Notes}

\begin{document}

\maketitle
\tableofcontents

We need to find some suitable norm for the vector space $(\Delta _n,\oplus;\mathbb{R},\odot)$, since the inner product $\left< p,q \right>=\sum\text{clr}_i(p)\text{clr}_i(q)$ forms a Hilbert space (with $L^{2}$ norm), which is isometric isomorphism to its dual space, but $(\Delta _n,\oplus;\mathbb{R},\odot) \not\cong(\mathbb{R},+,\cdot)$.


\chapter{Complete Analysis of Norms for Aitchison Geometry}

\section{1. Introduction and Basic Definitions}

\subsection{1.1 The Probability Simplex}

The \textbf{probability simplex} $\Delta_n$ is the set of all probability distributions over $n+1$ outcomes:
\[
\Delta_n = \left\{(p_0, p_1, \ldots, p_n) : p_i \geq 0 \text{ for all } i, \text{ and } \sum_{i=0}^n p_i = 1\right\}
\]

For example, if $n=2$, then $\Delta_2$ consists of all probability distributions over 3 outcomes, like $(0.2, 0.3, 0.5)$ or $(1/3, 1/3, 1/3)$.

\subsection{1.2 The Aitchison Geometry}

In ordinary arithmetic, we add numbers and multiply by scalars. But probability distributions live on a constrained space (they must sum to 1), so we need special operations. The \textbf{Aitchison geometry} provides a way to treat the simplex as a vector space with operations that respect this constraint.

\textbf{Vector Addition} $\oplus$:
\[
p \oplus q = \left[ \frac{p_i q_i}{\sum_{j=0}^n p_j q_j} \right]_{i=0}^n
\]

\textbf{Intuition}: We multiply corresponding components and then renormalize to get a probability distribution.

\textbf{Example}: Let $p = (0.2, 0.3, 0.5)$ and $q = (0.1, 0.6, 0.3)$. Then:

\begin{itemize}
	\item First multiply: $(0.2 \times 0.1, 0.3 \times 0.6, 0.5 \times 0.3) = (0.02, 0.18, 0.15)$
	\item Sum: $0.02 + 0.18 + 0.15 = 0.35$
	\item Normalize: $p \oplus q = (0.02/0.35, 0.18/0.35, 0.15/0.35) \approx (0.057, 0.514, 0.429)$
\end{itemize}

\textbf{Scalar Multiplication} $\odot$:
\[
c \odot p = \left[ \frac{(p_i)^c}{\sum_{j=0}^n (p_j)^c} \right]_{i=0}^n
\]

\textbf{Intuition}: We raise each component to the power $c$ and renormalize.

\textbf{Example}: Let $p = (0.2, 0.3, 0.5)$ and $c = 2$. Then:

\begin{itemize}
	\item First power: $(0.2^2, 0.3^2, 0.5^2) = (0.04, 0.09, 0.25)$
	\item Sum: $0.04 + 0.09 + 0.25 = 0.38$
	\item Normalize: $2 \odot p = (0.04/0.38, 0.09/0.38, 0.25/0.38) \approx (0.105, 0.237, 0.658)$
\end{itemize}

\textbf{Zero Element}: $\mathbf{o} = \left(\frac{1}{n+1}, \ldots, \frac{1}{n+1}\right)$ (the uniform distribution)

\subsection{1.3 The Centered Log-Ratio (clr) Transform}

The \textbf{clr transform} is the key to understanding this geometry. It maps probability distributions to vectors in Euclidean space:
\[
\text{clr}_i(p) = \log\frac{p_i}{g(p)} \quad \text{where } g(p) = \left(\prod_{j=0}^n p_j\right)^{1/(n+1)}
\]

Here $g(p)$ is the \textbf{geometric mean} of the components of $p$.

\textbf{Example}: For $p = (0.2, 0.3, 0.5)$:

\begin{itemize}
	\item Geometric mean: $g(p) = (0.2 \times 0.3 \times 0.5)^{1/3} = 0.03^{1/3} \approx 0.311$
	\item clr transform: $\text{clr}(p) = (\log(0.2/0.311), \log(0.3/0.311), \log(0.5/0.311)) \approx (-0.442, -0.036, 0.475)$
\end{itemize}

\textbf{Key Property}: The clr transform linearizes the Aitchison operations:

\begin{itemize}
	\item $\text{clr}(p \oplus q) = \text{clr}(p) + \text{clr}(q)$ (vector addition becomes ordinary addition)
	\item $\text{clr}(c \odot p) = c \cdot \text{clr}(p)$ (scalar multiplication becomes ordinary multiplication)
\end{itemize}

\textbf{Important}: The clr vectors always sum to zero: $\sum_{i=0}^n \text{clr}_i(p) = 0$.

\section{2. What is a Norm?}

A \textbf{norm} on a vector space is a function $\|\cdot\|$ that measures the "size" or "length" of vectors. It must satisfy three properties:

\begin{enumerate}
	\item \textbf{Positive definiteness}: $\|v\| \geq 0$ for all vectors $v$, and $\|v\| = 0$ if and only if $v$ is the zero vector
	\item \textbf{Homogeneity}: $\|\lambda v\| = |\lambda| \|v\|$ for any scalar $\lambda$ and vector $v$
	\item \textbf{Triangle inequality}: $\|u + v\| \leq \|u\| + \|v\|$ for any vectors $u, v$
\end{enumerate}

\section{3. Valid Norms on the Aitchison Space}

\subsection{3.1 The \texorpdfstring{$L^p$}{L^p} Norms (Complete Analysis)}

\subsubsection{Definition}

For $1 \leq p < \infty$:
\[
\|q\|_p = \left( \sum_{i=0}^n |\text{clr}_i(q)|^p \right)^{1/p}
\]

For $p = \infty$:
\[
\|q\|_\infty = \max_{i=0,\ldots,n} |\text{clr}_i(q)|
\]

\subsubsection{Complete Proof that \texorpdfstring{$\|\cdot\|_p$}{|cdot|_p} is a Norm}

\textbf{Proof of Positive Definiteness}:

\begin{itemize}
	\item Step 1: Since $|\text{clr}_i(q)|^p \geq 0$ for all $i$, we have $\|q\|_p \geq 0$.
	\item Step 2: If $\|q\|_p = 0$, then $\sum_{i=0}^n |\text{clr}_i(q)|^p = 0$.
	\item Step 3: Since each term is non-negative, this means $|\text{clr}_i(q)|^p = 0$ for all $i$.
	\item Step 4: Therefore $\text{clr}_i(q) = 0$ for all $i$.
	\item Step 5: This means $\log\frac{q_i}{g(q)} = 0$, so $q_i = g(q)$ for all $i$.
	\item Step 6: Since $\sum q_i = 1$ and all $q_i$ are equal, we must have $q_i = \frac{1}{n+1}$ for all $i$.
	\item Step 7: Therefore $q = \mathbf{o}$ (the zero element).
	\item Conversely, if $q = \mathbf{o}$, then $\text{clr}_i(\mathbf{o}) = \log\frac{1/(n+1)}{1/(n+1)} = 0$ for all $i$, so $\|\mathbf{o}\|_p = 0$.
\end{itemize}

\textbf{Proof of Homogeneity}:

\begin{itemize}
	\item Step 1: We need to show $\|\lambda \odot q\|_p = |\lambda| \|q\|_p$.
	\item Step 2: First, let's find $\text{clr}_i(\lambda \odot q)$.
	\item Step 3: $(\lambda \odot q)_i = \frac{(q_i)^\lambda}{\sum_{j=0}^n (q_j)^\lambda}$
	\item Step 4: The geometric mean of $\lambda \odot q$ is:
\[
g(\lambda \odot q) = \left(\prod_{i=0}^n \frac{(q_i)^\lambda}{\sum_{j=0}^n (q_j)^\lambda}\right)^{1/(n+1)} = \frac{\left(\prod_{i=0}^n (q_i)^\lambda\right)^{1/(n+1)}}{\sum_{j=0}^n (q_j)^\lambda}
\]
	\item Step 5: Therefore:
\[
\text{clr}_i(\lambda \odot q) = \log\frac{(\lambda \odot q)_i}{g(\lambda \odot q)} = \log\frac{(q_i)^\lambda}{\left(\prod_{j=0}^n (q_j)^\lambda\right)^{1/(n+1)}} = \lambda \log\frac{q_i}{g(q)} = \lambda \cdot \text{clr}_i(q)
\]
	\item Step 6: Now:
\[
\|\lambda \odot q\|_p = \left(\sum_{i=0}^n |\lambda \cdot \text{clr}_i(q)|^p\right)^{1/p} = \left(|\lambda|^p \sum_{i=0}^n |\text{clr}_i(q)|^p\right)^{1/p} = |\lambda| \|q\|_p
\]
\end{itemize}

\textbf{Proof of Triangle Inequality}:

\begin{itemize}
	\item Step 1: We need to show $\|p \oplus q\|_p \leq \|p\|_p + \|q\|_p$.
	\item Step 2: Using the linearity of clr: $\text{clr}(p \oplus q) = \text{clr}(p) + \text{clr}(q)$.
	\item Step 3: Therefore:
\[
\|p \oplus q\|_p = \left(\sum_{i=0}^n |\text{clr}_i(p) + \text{clr}_i(q)|^p\right)^{1/p}
\]
	\item Step 4: By the Minkowski inequality (which is the triangle inequality for $L^p$ spaces in $\mathbb{R}^{n+1}$):
\[
\left(\sum_{i=0}^n |a_i + b_i|^p\right)^{1/p} \leq \left(\sum_{i=0}^n |a_i|^p\right)^{1/p} + \left(\sum_{i=0}^n |b_i|^p\right)^{1/p}
\]
	\item Step 5: Applying this with $a_i = \text{clr}_i(p)$ and $b_i = \text{clr}_i(q)$:
\[
\|p \oplus q\|_p \leq \|p\|_p + \|q\|_p
\]
\end{itemize}

\textbf{Special Cases}:

\begin{itemize}
	\item $p = 1$: The "taxicab" norm, $\|q\|_1 = \sum_{i=0}^n |\text{clr}_i(q)|$
	\item $p = 2$: The Euclidean norm, $\|q\|_2 = \sqrt{\sum_{i=0}^n [\text{clr}_i(q)]^2}$ (induces the Aitchison inner product)
	\item $p = \infty$: The maximum norm, $\|q\|_\infty = \max_i |\text{clr}_i(q)|$
\end{itemize}

\subsection{3.2 The Dual Norm}

For any norm $\|\cdot\|$ on a vector space $V$, the \textbf{dual norm} $\|\cdot\|_*$ on the dual space $V^*$ is defined by:
\[
\|f\|_* = \sup_{\|v\| \leq 1} |\langle f, v \rangle|
\]

where $\langle f, v \rangle$ is the pairing between $f \in V^*$ and $v \in V$.

\textbf{For $L^p$ norms}: The dual of $\|\cdot\|_p$ is $\|\cdot\|_q$ where $\frac{1}{p} + \frac{1}{q} = 1$.

\textbf{Example}:

\begin{itemize}
	\item The dual of $\|\cdot\|_2$ is $\|\cdot\|_2$ itself (self-dual)
	\item The dual of $\|\cdot\|_1$ is $\|\cdot\|_\infty$
	\item The dual of $\|\cdot\|_\infty$ is $\|\cdot\|_1$
\end{itemize}

\subsection{3.3 Weighted \texorpdfstring{$L^p$}{L^p} Norms}

\subsubsection{Definition}

For fixed weights $w_i > 0$:
\[
\|q\|_{p,w} = \left( \sum_{i=0}^n w_i |\text{clr}_i(q)|^p \right)^{1/p}
\]

\textbf{Intuition}: Different components can have different importance. For instance, if we're more interested in changes to certain probabilities, we can weight them more heavily.

The proof that this is a norm follows the same steps as for the unweighted case, with the weights carried through.

\textbf{Dual norm}: $\|f\|_{p,w,*} = \left( \sum_{i=0}^n w_i^{-q/p} |f_i|^q \right)^{1/q}$ where $\frac{1}{p} + \frac{1}{q} = 1$.

\subsection{3.3 Weighted \texorpdfstring{$L^p$}{L^p} Norms (Complete Analysis)}

\subsubsection{Definition and Motivation}

For fixed weights $w = (w_0, w_1, \ldots, w_n)$ with $w_i > 0$:
\[
\|q\|_{p,w} = \left( \sum_{i=0}^n w_i |\text{clr}_i(q)|^p \right)^{1/p}
\]

\textbf{Intuition}: In many applications, different probability components have different significance:

\begin{itemize}
	\item In medical diagnosis, certain conditions might be more critical
	\item In financial portfolios, some assets might be more important
	\item In ecological systems, certain species might be keystone species
\end{itemize}

The weights $w_i$ allow us to emphasize these differences.

\subsubsection{Complete Proof that \texorpdfstring{$\|\cdot\|_{p,w}$}{|cdot|_p,w} is a Norm}

\textbf{Proof of Positive Definiteness}:

\begin{itemize}
	\item Step 1: Since $w_i > 0$ and $|\text{clr}_i(q)|^p \geq 0$, we have $w_i|\text{clr}_i(q)|^p \geq 0$ for all $i$.
	\item Step 2: Therefore $\|q\|_{p,w} \geq 0$.
	\item Step 3: If $\|q\|_{p,w} = 0$, then $\sum_{i=0}^n w_i|\text{clr}_i(q)|^p = 0$.
	\item Step 4: Since each term $w_i|\text{clr}_i(q)|^p \geq 0$, we must have $w_i|\text{clr}_i(q)|^p = 0$ for all $i$.
	\item Step 5: Since $w_i > 0$, this implies $|\text{clr}_i(q)|^p = 0$, hence $\text{clr}_i(q) = 0$ for all $i$.
	\item Step 6: As shown before, this means $q = \mathbf{o}$.
	\item Conversely, if $q = \mathbf{o}$, then $\text{clr}_i(\mathbf{o}) = 0$ for all $i$, so $\|\mathbf{o}\|_{p,w} = 0$.
\end{itemize}

\textbf{Proof of Homogeneity}:

\begin{itemize}
	\item Step 1: We know that $\text{clr}_i(\lambda \odot q) = \lambda \cdot \text{clr}_i(q)$.
	\item Step 2: Therefore:
\[
\|\lambda \odot q\|_{p,w} = \left(\sum_{i=0}^n w_i |\lambda \cdot \text{clr}_i(q)|^p\right)^{1/p}
\]
	\item Step 3: Factor out $|\lambda|^p$:
\[
= \left(|\lambda|^p \sum_{i=0}^n w_i |\text{clr}_i(q)|^p\right)^{1/p}
\]
	\item Step 4: Take the $p$-th root:
\[
= |\lambda| \left(\sum_{i=0}^n w_i |\text{clr}_i(q)|^p\right)^{1/p} = |\lambda| \|q\|_{p,w}
\]
\end{itemize}

\textbf{Proof of Triangle Inequality}:

\begin{itemize}
	\item Step 1: We have $\text{clr}(p \oplus q) = \text{clr}(p) + \text{clr}(q)$.
	\item Step 2: By the weighted Minkowski inequality (generalization of the standard Minkowski inequality):
\[
\left(\sum_{i=0}^n w_i |a_i + b_i|^p\right)^{1/p} \leq \left(\sum_{i=0}^n w_i |a_i|^p\right)^{1/p} + \left(\sum_{i=0}^n w_i |b_i|^p\right)^{1/p}
\]
	\item Step 3: Apply this with $a_i = \text{clr}_i(p)$ and $b_i = \text{clr}_i(q)$:
\[
\|p \oplus q\|_{p,w} \leq \|p\|_{p,w} + \|q\|_{p,w}
\]
\end{itemize}

\subsubsection{Derivation of the Dual Norm}

\textbf{Theorem}: The dual norm of $\|\cdot\|_{p,w}$ is $\|f\|_{p,w,*} = \left( \sum_{i=0}^n w_i^{-q/p} |f_i|^q \right)^{1/q}$ where $\frac{1}{p} + \frac{1}{q} = 1$.

\textbf{Proof}:

\begin{itemize}
	\item Step 1: By definition, $\|f\|_{p,w,*} = \sup_{\|x\|_{p,w} \leq 1} |\langle f, x \rangle|$.
	\item Step 2: Using Hölder's inequality with weights, for any $x$ with $\|x\|_{p,w} \leq 1$:
\[
|\langle f, x \rangle| = \left|\sum_{i=0}^n f_i x_i\right| \leq \sum_{i=0}^n |f_i| |x_i|
\]
	\item Step 3: Rewrite using the weights:
\[
= \sum_{i=0}^n (w_i^{1/p} |x_i|) \cdot (w_i^{-1/p} |f_i|)
\]
	\item Step 4: Apply Hölder's inequality:
\[
\leq \left(\sum_{i=0}^n w_i |x_i|^p\right)^{1/p} \left(\sum_{i=0}^n w_i^{-q/p} |f_i|^q\right)^{1/q}
\]
	\item Step 5: Since $\|x\|_{p,w} \leq 1$, the first factor is at most 1:
\[
|\langle f, x \rangle| \leq \left(\sum_{i=0}^n w_i^{-q/p} |f_i|^q\right)^{1/q}
\]
	\item Step 6: This bound is achieved when $x_i = \text{sign}(f_i) \cdot w_i^{-1/p} |f_i|^{q-1} / \|f\|_{p,w,*}^{q-1}$.
	\item Step 7: Therefore $\|f\|_{p,w,*} = \left(\sum_{i=0}^n w_i^{-q/p} |f_i|^q\right)^{1/q}$.
\end{itemize}

\textbf{Key Insight}: The weights appear with negative exponents in the dual norm, reflecting an inverse relationship between primal and dual weights.

\subsection{3.4 Orlicz Norms (Advanced Topic)}

\subsubsection{What is a Young Function?}

A \textbf{Young function} $\Phi: \mathbb{R} \to [0, \infty)$ is a function that:

\begin{enumerate}
	\item Is convex (curves upward)
	\item Is even: $\Phi(-x) = \Phi(x)$
	\item Satisfies $\Phi(0) = 0$
	\item Goes to infinity: $\Phi(x) \to \infty$ as $|x| \to \infty$
\end{enumerate}

\textbf{Examples}:

\begin{itemize}
	\item $\Phi(x) = |x|^p$ for $p \geq 1$ (gives $L^p$ norms)
	\item $\Phi(x) = e^{|x|} - 1$ (gives exponential growth)
	\item $\Phi(x) = |x| \log(1 + |x|)$ (between linear and quadratic)
\end{itemize}

\subsubsection{Definition of Orlicz Norm}
\[
\|q\|_\Phi = \inf \left\{ k > 0 : \sum_{i=0}^n \Phi\left(\frac{|\text{clr}_i(q)|}{k}\right) \leq 1 \right\}
\]

\textbf{Intuition}: We scale down the clr vector by a factor $k$ until the sum of $\Phi$ applied to each scaled component is at most 1. The smallest such $k$ is the norm.

\subsection{3.4 Orlicz Norms (Complete Analysis)}

\subsubsection{What is a Young Function?}

A \textbf{Young function} $\Phi: \mathbb{R} \to [0, \infty)$ is a function satisfying:

\begin{enumerate}
	\item \textbf{Convexity}: $\Phi(\alpha x + (1-\alpha)y) \leq \alpha\Phi(x) + (1-\alpha)\Phi(y)$ for $\alpha \in [0,1]$
	\item \textbf{Even}: $\Phi(-x) = \Phi(x)$
	\item \textbf{Zero at origin}: $\Phi(0) = 0$
	\item \textbf{Growth condition}: $\lim_{|x| \to \infty} \Phi(x) = \infty$
	\item \textbf{Superlinearity}: $\lim_{|x| \to \infty} \frac{\Phi(x)}{|x|} = \infty$
\end{enumerate}

\textbf{Examples with Properties}:

\begin{itemize}
	\item $\Phi(x) = |x|^p$ for $p > 1$: Grows polynomially
	\item $\Phi(x) = e^{|x|} - 1$: Grows exponentially fast
	\item $\Phi(x) = |x| \log(1 + |x|)$: Between linear and quadratic growth
	\item $\Phi(x) = \begin{cases} 0 & \text{if } |x| \leq 1 \\ \infty & \text{if } |x| > 1 \end{cases}$: Gives the $\ell^\infty$ norm
\end{itemize}

\subsubsection{Definition and Alternative Characterizations}

The \textbf{Orlicz norm} can be defined in three equivalent ways:

\begin{enumerate}
	\item \textbf{Gauge form} (what we use):
\[
\|q\|_\Phi = \inf \left\{ k > 0 : \sum_{i=0}^n \Phi\left(\frac{|\text{clr}_i(q)|}{k}\right) \leq 1 \right\}
\]
	\item \textbf{Amemiya form}:
\[
\|q\|_\Phi = \inf_{k > 0} \left\{ k + \frac{1}{k} \sum_{i=0}^n \Phi(|\text{clr}_i(q)|) \right\}
\]
	\item \textbf{Luxemburg form}:
\[
\|q\|_\Phi = \inf \left\{ \sum_{i=0}^n \Phi\left(\frac{|\text{clr}_i(q)|}{\lambda}\right) \leq 1 \right\}
\]
\end{enumerate}

\subsubsection{Complete Proof that \texorpdfstring{$\|\cdot\|_\Phi$}{|cdot|_Phi} is a Norm}

\textbf{Proof of Positive Definiteness}:

\begin{itemize}
	\item Step 1: Clearly $\|q\|_\Phi \geq 0$ since it's an infimum over positive numbers.
	\item Step 2: If $q = \mathbf{o}$, then $\text{clr}_i(q) = 0$ for all $i$.
	\item Step 3: Thus $\Phi(0/k) = \Phi(0) = 0$ for any $k > 0$.
	\item Step 4: So $\sum \Phi(|\text{clr}_i(\mathbf{o})|/k) = 0 \leq 1$ for any $k > 0$.
	\item Step 5: Therefore $\|\mathbf{o}\|_\Phi = \inf\{k > 0\} = 0$.
	\item Step 6: Conversely, suppose $\|q\|_\Phi = 0$.
	\item Step 7: This means for any $\epsilon > 0$, there exists $k < \epsilon$ with $\sum \Phi(|\text{clr}_i(q)|/k) \leq 1$.
	\item Step 8: As $k \to 0$, if any $\text{clr}_i(q) \neq 0$, then $\Phi(|\text{clr}_i(q)|/k) \to \infty$ by the growth condition.
	\item Step 9: This would make the sum exceed 1, a contradiction.
	\item Step 10: Therefore all $\text{clr}_i(q) = 0$, so $q = \mathbf{o}$.
\end{itemize}

\textbf{Proof of Homogeneity}:

\begin{itemize}
	\item Step 1: We need to show $\|\lambda \odot q\|_\Phi = |\lambda| \|q\|_\Phi$.
	\item Step 2: Since $\text{clr}_i(\lambda \odot q) = \lambda \cdot \text{clr}_i(q)$:
\[
\|\lambda \odot q\|_\Phi = \inf \left\{ k > 0 : \sum_{i=0}^n \Phi\left(\frac{|\lambda \cdot \text{clr}_i(q)|}{k}\right) \leq 1 \right\}
\]
	\item Step 3: Substitute $k = |\lambda| k'$:
\[
= \inf \left\{ |\lambda|k' > 0 : \sum_{i=0}^n \Phi\left(\frac{|\text{clr}_i(q)|}{k'}\right) \leq 1 \right\}
\]
	\item Step 4: Factor out $|\lambda|$:
\[
= |\lambda| \inf \left\{ k' > 0 : \sum_{i=0}^n \Phi\left(\frac{|\text{clr}_i(q)|}{k'}\right) \leq 1 \right\} = |\lambda| \|q\|_\Phi
\]
\end{itemize}

\textbf{Proof of Triangle Inequality}:

\begin{itemize}
	\item Step 1: We need to show $\|p \oplus q\|_\Phi \leq \|p\|_\Phi + \|q\|_\Phi$.
	\item Step 2: Let $k_1 = \|p\|_\Phi + \epsilon$ and $k_2 = \|q\|_\Phi + \epsilon$ for small $\epsilon > 0$.
	\item Step 3: By definition, $\sum \Phi(|\text{clr}_i(p)|/k_1) \leq 1$ and $\sum \Phi(|\text{clr}_i(q)|/k_2) \leq 1$.
	\item Step 4: For $\theta = k_1/(k_1 + k_2)$, by convexity of $\Phi$:
\[
\Phi\left(\frac{|\text{clr}_i(p) + \text{clr}_i(q)|}{k_1 + k_2}\right) \leq \theta \Phi\left(\frac{|\text{clr}_i(p)|}{k_1}\right) + (1-\theta) \Phi\left(\frac{|\text{clr}_i(q)|}{k_2}\right)
\]
	\item Step 5: Sum over $i$:
\[
\sum_{i=0}^n \Phi\left(\frac{|\text{clr}_i(p \oplus q)|}{k_1 + k_2}\right) \leq \theta \cdot 1 + (1-\theta) \cdot 1 = 1
\]
	\item Step 6: Therefore $\|p \oplus q\|_\Phi \leq k_1 + k_2 = \|p\|_\Phi + \|q\|_\Phi + 2\epsilon$.
	\item Step 7: Since this holds for any $\epsilon > 0$, we get $\|p \oplus q\|_\Phi \leq \|p\|_\Phi + \|q\|_\Phi$.
\end{itemize}

\subsubsection{The Conjugate Young Function and Dual Norm}

\textbf{Definition}: For a Young function $\Phi$, its \textbf{conjugate} (or complementary) Young function is:
\[
\Psi(y) = \sup_{x \geq 0} \{xy - \Phi(x)\}
\]

\textbf{Examples of Conjugate Pairs}:

\begin{enumerate}
	\item If $\Phi(x) = |x|^p/p$ with $p > 1$, then $\Psi(y) = |y|^q/q$ where $1/p + 1/q = 1$
	\item If $\Phi(x) = e^{|x|} - 1$, then $\Psi(y) = |y|\log|y| - |y| + 1$ for $|y| > 1$
	\item If $\Phi(x) = |x|\log(1 + |x|)$, then $\Psi(y) \approx e^{|y|} - 1$
\end{enumerate}

\textbf{Theorem}: The dual norm of $\|\cdot\|_\Phi$ is $\|\cdot\|_\Psi$.

\textbf{Proof Sketch}:

\begin{itemize}
	\item Step 1: By the Young inequality: $xy \leq \Phi(x) + \Psi(y)$ for all $x, y$.
	\item Step 2: For any $f$ and $x$ with $\sum \Phi(|x_i|) \leq 1$:
\[
|\langle f, x \rangle| \leq \sum_{i=0}^n |f_i||x_i| \leq \sum_{i=0}^n [\Phi(|x_i|) + \Psi(|f_i|)] \leq 1 + \sum_{i=0}^n \Psi(|f_i|)
\]
	\item Step 3: This gives $\|f\|_* \leq \|f\|_\Psi$.
	\item Step 4: The reverse inequality follows from the fact that equality in Young's inequality can be achieved.
\end{itemize}

\textbf{Geometric Interpretation}:

\begin{itemize}
	\item The primal Orlicz norm measures size according to the growth rate of $\Phi$
	\item The dual norm measures size according to the complementary growth rate $\Psi$
	\item Fast growth in the primal (e.g., exponential) corresponds to slow growth in the dual (e.g., logarithmic)
\end{itemize}

\section{4. Connection to Information Geometry}

\subsection{4.1 What are Exponential and Mixture Families?}

\subsubsection{Exponential Families (e-families)}

An \textbf{exponential family} is a family of probability distributions whose log-probabilities are linear in some parameters.

\textbf{Example}: Consider categorical distributions on $\{0, 1, \ldots, n\}$. We can write:
\[
\log p_i = \theta_i - \psi(\theta)
\]
where $\theta = (\theta_0, \ldots, \theta_n)$ are the \textbf{natural parameters} and $\psi(\theta)$ is a normalization factor ensuring $\sum p_i = 1$.

\textbf{Key insight}: The clr transform gives us exactly these natural parameters!
\[
\text{clr}_i(p) = \log p_i - \frac{1}{n+1}\sum_{j=0}^n \log p_j \approx \theta_i - \text{constant}
\]

\subsubsection{Mixture Families (m-families)}

A \textbf{mixture family} consists of distributions formed by taking weighted averages (convex combinations) of other distributions:
\[
p = \sum_{j} w_j q^{(j)} \quad \text{where } w_j \geq 0, \sum w_j = 1
\]

\textbf{Example}: If we have three "pure" strategies with distributions $q^{(1)}, q^{(2)}, q^{(3)}$, then any mixture like $0.5 q^{(1)} + 0.3 q^{(2)} + 0.2 q^{(3)}$ is in the mixture family.

\subsection{4.2 The Duality}

In information geometry, these two families are \textbf{dual} to each other:

\begin{itemize}
	\item \textbf{e-flat coordinates}: The natural parameters $\theta$ (what clr gives us)
	\item \textbf{m-flat coordinates}: The expectation parameters $\eta$ (mixture weights)
\end{itemize}

The \textbf{Fisher information metric} provides the Riemannian structure that makes these two coordinate systems dual.

\subsection{4.3 The Deep Connection: How Norm Duality Reflects e/m Duality}

\subsubsection{4.3.1 The Fundamental Theorem}

\textbf{Theorem}: The duality between norms $\|\cdot\|_p$ and $\|\cdot\|_q$ (where $\frac{1}{p} + \frac{1}{q} = 1$) on the clr-transformed space is geometrically equivalent to the duality between e-flat and m-flat coordinate systems in information geometry.

\textbf{Proof Strategy}:

\begin{enumerate}
	\item Show that the clr transform maps the simplex to the e-flat coordinate system
	\item Demonstrate that the dual space corresponds to the m-flat coordinate system
	\item Prove that the Legendre transform relating e/m coordinates corresponds to the norm duality
\end{enumerate}

\subsubsection{4.3.2 The clr Transform as e-flat Coordinates}

\textbf{Detailed Construction}:

Consider the exponential family of categorical distributions:
\[
p_i(\theta) = \frac{\exp(\theta_i)}{\sum_{j=0}^n \exp(\theta_j)}
\]

The \textbf{log-partition function} is:
\[
\psi(\theta) = \log\left(\sum_{j=0}^n \exp(\theta_j)\right)
\]

Taking the logarithm:
\[
\log p_i(\theta) = \theta_i - \psi(\theta)
\]

\textbf{Connection to clr}:
The clr transform gives us:
\[
\text{clr}_i(p) = \log p_i - \frac{1}{n+1}\sum_{j=0}^n \log p_j
\]

This is precisely the natural parameter $\theta_i$ up to a constant shift! More precisely:
\[
\text{clr}_i(p) = \theta_i - \frac{1}{n+1}\sum_{j=0}^n \theta_j
\]

\textbf{Key Insight}: The clr space is exactly the e-flat coordinate system, and our norms $\|\cdot\|_p$ on clr space are measuring distances in the natural parameter space.

\subsubsection{4.3.3 The Dual Space as m-flat Coordinates}

\textbf{The Expectation Parameters}:
For the exponential family $p_i(\theta) = \frac{\exp(\theta_i)}{\sum_j \exp(\theta_j)}$, the \textbf{expectation parameters} are:
\[
\eta_i = \mathbb{E}[T_i] = \frac{\partial \psi(\theta)}{\partial \theta_i} = \frac{\exp(\theta_i)}{\sum_j \exp(\theta_j)} = p_i
\]

\textbf{Crucial Observation}: The expectation parameters $\eta$ are exactly the probability values $p$ themselves!

\textbf{The Dual Pairing}:
The duality between e-flat and m-flat coordinates is expressed through the \textbf{dual pairing}:
\[
\langle \theta, \eta \rangle = \sum_{i=0}^n \theta_i \eta_i = \sum_{i=0}^n \theta_i p_i
\]

In clr coordinates, this becomes:
\[
\langle \text{clr}(p), p \rangle = \sum_{i=0}^n \text{clr}_i(p) \cdot p_i
\]

\subsubsection{4.3.4 The Fisher Information Metric and \texorpdfstring{$L^2$}{L^2} Norm}

\textbf{The Fisher Information Matrix}:
For the categorical exponential family, the Fisher information matrix is:
\[
G_{ij}(\theta) = \frac{\partial^2 \psi(\theta)}{\partial \theta_i \partial \theta_j} = \delta_{ij} p_i - p_i p_j
\]

where $\delta_{ij}$ is the Kronecker delta.

\textbf{In clr coordinates}, this becomes:
\[
G_{ij}^{\text{clr}} = \delta_{ij} p_i + \frac{1}{n+1} \sum_{k=0}^n p_k
\]

\textbf{The Induced Metric}:
The Fisher information metric in clr coordinates is:
\[
\langle u, v \rangle_{\text{Fisher}} = \sum_{i,j=0}^n u_i G_{ij}^{\text{clr}} v_j
\]

\textbf{Remarkable Fact}: When we restrict to the subspace where $\sum u_i = 0$ (which is exactly the clr constraint), this reduces to:
\[
\langle u, v \rangle_{\text{Fisher}} = \sum_{i=0}^n u_i v_i
\]

This is precisely the standard Euclidean inner product, which induces the $L^2$ norm!

\textbf{Theorem}: The Aitchison $L^2$ norm $\|q\|_2 = \sqrt{\sum_{i=0}^n [\text{clr}_i(q)]^2}$ is the natural Riemannian distance induced by the Fisher information metric on the probability simplex.

\subsubsection{4.3.5 Norm Duality and Coordinate System Duality}

\textbf{The Legendre Transform}:
The relationship between e-flat and m-flat coordinates is given by the \textbf{Legendre transform}:
\[
\psi^*(\eta) = \sup_\theta \{\langle \theta, \eta \rangle - \psi(\theta)\}
\]

This is the \textbf{convex conjugate} of the log-partition function.

\textbf{Connection to Norm Duality}:
The duality between $\|\cdot\|_p$ and $\|\cdot\|_q$ is also expressed through a supremum:
\[
\|f\|_q = \sup_{\|x\|_p \leq 1} |\langle f, x \rangle|
\]

\textbf{Theorem}: The norm duality $\|\cdot\|_p \leftrightarrow \|\cdot\|_q$ is the functional analytic manifestation of the Legendre transform duality between e-flat and m-flat coordinates.

\textbf{Proof Sketch}:

\begin{enumerate}
	\item The constraint $\|x\|_p \leq 1$ in clr space corresponds to a constraint on the natural parameters $\theta$
	\item The supremum $\sup_{\|x\|_p \leq 1} |\langle f, x \rangle|$ corresponds to the Legendre transform $\sup_\theta \{\langle \theta, \eta \rangle - \psi(\theta)\}$
	\item The dual norm $\|f\|_q$ measures distance in the expectation parameter space (m-flat coordinates)
\end{enumerate}

\subsubsection{4.3.6 Specific Examples of the Duality}

\textbf{Example 1: $L^1$ and $L^\infty$ Duality}

\begin{itemize}
	\item \textbf{$L^1$ norm on clr space}: $\|q\|_1 = \sum_{i=0}^n |\text{clr}_i(q)|$
	\begin{itemize}
		\item \textbf{Geometric meaning}: Measures "taxicab distance" in natural parameter space
		\item \textbf{Information theoretic meaning}: Related to the total variation distance
	\end{itemize}
	\item \textbf{$L^\infty$ norm on clr space}: $\|q\|_\infty = \max_i |\text{clr}_i(q)|$
	\begin{itemize}
		\item \textbf{Geometric meaning}: Measures "maximum deviation" in natural parameter space
		\item \textbf{Information theoretic meaning}: Related to the maximum likelihood principle
	\end{itemize}
\end{itemize}

\textbf{The Duality}:
\[
\|f\|_\infty = \sup_{\|x\|_1 \leq 1} |\langle f, x \rangle|
\]

This reflects the fact that extreme points of the $L^1$ ball correspond to coordinate directions, which are exactly the pure strategies in the mixture family.

\textbf{Example 2: Weighted Norms and Importance Sampling}

\begin{itemize}
	\item \textbf{Weighted $L^p$ norm}: $\|q\|_{p,w} = \left( \sum_{i=0}^n w_i |\text{clr}_i(q)|^p \right)^{1/p}$
	\begin{itemize}
		\item \textbf{Geometric meaning}: Different importance weights $w_i$ for different outcomes
		\item \textbf{Information theoretic meaning}: Corresponds to importance sampling with weights $w_i$
	\end{itemize}
	\item \textbf{Dual weighted norm}: $\|f\|_{q,w^{-1}} = \left( \sum_{i=0}^n w_i^{-q/p} |f_i|^q \right)^{1/q}$
	\begin{itemize}
		\item \textbf{Geometric meaning}: Inverse weighting in the expectation parameter space
		\item \textbf{Information theoretic meaning}: Reflects the bias-variance tradeoff in importance sampling
	\end{itemize}
\end{itemize}

\subsubsection{4.3.7 Orlicz Norms and Divergences}

\textbf{The Connection to f-Divergences}:
Orlicz norms with Young function $\Phi$ are intimately connected to \textbf{f-divergences} in information theory.

\textbf{Definition}: The f-divergence between distributions $p$ and $q$ is:
\[
D_f(p\|q) = \sum_{i=0}^n q_i f\left(\frac{p_i}{q_i}\right)
\]

\textbf{Theorem}: The Orlicz norm $\|p\|_\Phi$ is related to the f-divergence $D_\Phi(p\|\mathbf{o})$ where $\mathbf{o}$ is the uniform distribution.

\textbf{Proof}:
\[
\|p\|_\Phi = \inf \left\{ k > 0 : \sum_{i=0}^n \Phi\left(\frac{|\text{clr}_i(p)|}{k}\right) \leq 1 \right\}
\]

Setting $f(x) = \Phi(\log x)$ and using the fact that $\text{clr}_i(p) = \log p_i - \log g(p)$:
\[
\|p\|_\Phi \approx \sqrt{D_f(p\|\mathbf{o})}
\]

\textbf{Examples}:

\begin{enumerate}
	\item \textbf{$\Phi(x) = |x|^2/2$}: Gives the $L^2$ norm, corresponds to the $\chi^2$ divergence
	\item \textbf{$\Phi(x) = e^{|x|} - 1$}: Gives exponential Orlicz norm, corresponds to the exponential divergence
	\item \textbf{$\Phi(x) = |x|\log(1+|x|)$}: Gives logarithmic Orlicz norm, related to the Hellinger distance
\end{enumerate}

\textbf{The Duality}: The conjugate Young function $\Psi$ corresponds to the conjugate f-divergence, reflecting the duality between primal and dual divergences in information geometry.

\subsection{4.4 Summary of the Deep Connection}

The connection between suitable norms and e/m duality can be summarized as follows:

\begin{enumerate}
	\item \textbf{Structural Correspondence}:
	\begin{itemize}
		\item clr space ↔ e-flat coordinates (natural parameters)
		\item Dual of clr space ↔ m-flat coordinates (expectation parameters)
		\item Norm duality ↔ Legendre transform duality
	\end{itemize}
	\item \textbf{Metric Correspondence}:
	\begin{itemize}
		\item $L^2$ norm ↔ Fisher information metric
		\item $L^p$ norms ↔ $L^p$ Banach space structure on parameter spaces
		\item Orlicz norms ↔ f-divergences
	\end{itemize}
	\item \textbf{Functional Correspondence}:
	\begin{itemize}
		\item $\|\cdot\|_p$ measures distance in natural parameter space
		\item $\|\cdot\|_q$ (dual norm) measures distance in expectation parameter space
		\item Weighted norms ↔ Importance sampling and bias corrections
	\end{itemize}
	\item \textbf{Geometric Correspondence}:
	\begin{itemize}
		\item Convex hull in clr space ↔ m-flat submanifolds
		\item Affine subspaces in clr space ↔ e-flat submanifolds
		\item Dual norm balls ↔ Dual coordinate neighborhoods
	\end{itemize}
\end{enumerate}

This deep connection explains why the norms we identified are not just mathematically valid, but are the \textbf{natural} choices from the perspective of information geometry. They reflect the fundamental duality structure of statistical manifolds and provide the appropriate geometric framework for analyzing probability distributions.

\section{5. Why Other "Natural" Candidates Fail}

Several seemingly natural choices fail to be norms:

\subsection{5.1 Fisher Information Weighted Norm}
\[
\|q\|_F = \sqrt{\sum_{i=0}^n \frac{[\text{clr}_i(q)]^2}{q_i}} \quad \text{(FAILS!)}
\]

\textbf{Why it fails}: Under $\lambda \odot q$, the denominator $q_i$ transforms to $\frac{(q_i)^\lambda}{\sum (q_j)^\lambda}$, which doesn't scale properly with $\lambda$.

\subsection{5.2 Bregman-Exponential Norm}
\[
\|q\|_{Breg} = \sqrt{\sum_{i=0}^n [\text{clr}_i(q)]^2 \cdot \exp(\text{clr}_i(q))} \quad \text{(FAILS!)}
\]

\textbf{Why it fails}: The exponential term $\exp(\lambda \cdot \text{clr}_i(q)) \neq \lambda \cdot \exp(\text{clr}_i(q))$, breaking homogeneity.

\subsection{5.3 KL-Divergence Based Norm}
\[
\|q\|_{KL} = \sqrt{D_{KL}(q\|\mathbf{o}) + D_{KL}(\mathbf{o}\|q)} \quad \text{(FAILS!)}
\]

\textbf{Why it fails}: KL divergence doesn't scale linearly with the Aitchison scalar multiplication.

\textbf{Key Lesson}: Valid norms must work with the linear structure after clr transformation. Quantities that depend non-linearly on the original probability values typically fail.

\section{6. The Most Suitable Norms: Weighted \texorpdfstring{$L^p$}{L^p} and Orlicz Norms}

\subsection{6.1 Why the Standard \texorpdfstring{$L^2$}{L^2} Norm is Insufficient}

While the inner product $\langle p,q \rangle = \sum \text{clr}_i(p)\text{clr}_i(q)$ forms a Hilbert space, the standard $L^2$ norm has fundamental limitations in the Aitchison geometry:

\begin{enumerate}
	\item \textbf{Constraint Violation}: The clr constraint $\sum_{i=0}^n \text{clr}_i(p) = 0$ means we're working in a proper subspace of $\mathbb{R}^{n+1}$, not the full space
	\item \textbf{Isomorphism Problem}: $(\Delta_n,\oplus;\mathbb{R},\odot) \not\cong (\mathbb{R},+,\cdot)$ due to the geometric constraints
	\item \textbf{Uniform Weighting Issue}: The standard $L^2$ norm treats all probability components equally, ignoring their relative importance and natural variability
\end{enumerate}

\subsection{6.2 Weighted \texorpdfstring{$L^p$}{L^p} Norms: The Natural Generalization}

\subsubsection{6.2.1 Why Weighted \texorpdfstring{$L^p$}{L^p} Norms are More Suitable}

\textbf{Definition}:
\[
\|q\|_{p,w} = \left( \sum_{i=0}^n w_i |\text{clr}_i(q)|^p \right)^{1/p}
\]

\textbf{Key Advantages}:

\begin{enumerate}
	\item \textbf{Adaptive Scaling}: The weights $w_i$ can account for the natural variability of different probability components
	\item \textbf{Information-Theoretic Motivation}: Weights can reflect the Fisher information or other importance measures
	\item \textbf{Flexibility}: Different choices of $p$ and weights $w$ capture different aspects of the geometry
\end{enumerate}

\subsubsection{6.2.2 Detailed Interpretations of Weighted \texorpdfstring{$L^p$}{L^p} Norms}

\textbf{Case 1: $p = 1$ (Weighted $L^1$ Norm)}
\[
\|q\|_{1,w} = \sum_{i=0}^n w_i |\text{clr}_i(q)|
\]

\textbf{Geometric Interpretation}:

\begin{itemize}
	\item Measures weighted "taxicab distance" in the natural parameter space
	\item The unit ball is a weighted cross-polytope (hyperoctahedron)
	\item Promotes sparsity in the clr representation
\end{itemize}

\textbf{Information-Theoretic Interpretation}:

\begin{itemize}
	\item Related to weighted total variation distance
	\item Optimal for robust estimation when some components are more reliable
	\item Connection to $L^1$ regularization in high-dimensional statistics
\end{itemize}

\textbf{Practical Applications}:

\begin{itemize}
	\item \textbf{Ecological data}: Weight rare species more heavily to prevent extinction bias
	\item \textbf{Financial portfolios}: Weight volatile assets more to account for risk
	\item \textbf{Medical diagnosis}: Weight critical symptoms more heavily
\end{itemize}

\textbf{Example}: For ecological data where species $i$ has baseline abundance $\pi_i$, choose weights $w_i = 1/\pi_i$ to give equal importance to relative changes.

\textbf{Case 2: $p = \infty$ (Weighted $L^\infty$ Norm)}
\[
\|q\|_{\infty,w} = \max_{i=0,\ldots,n} w_i |\text{clr}_i(q)|
\]

\textbf{Geometric Interpretation}:

\begin{itemize}
	\item Measures the maximum weighted deviation in natural parameter space
	\item The unit ball is a weighted hypercube
	\item Ensures all weighted components stay within bounds
\end{itemize}

\textbf{Information-Theoretic Interpretation}:

\begin{itemize}
	\item Related to the maximum likelihood principle with weights
	\item Optimal for minimax estimation
	\item Connection to robust statistics and outlier detection
\end{itemize}

\textbf{Practical Applications}:

\begin{itemize}
	\item \textbf{Quality control}: Ensure no component exceeds its weighted tolerance
	\item \textbf{Risk management}: Control maximum weighted exposure
	\item \textbf{Fairness constraints}: Ensure equitable treatment across groups
\end{itemize}

\textbf{Case 3: $1 < p < \infty$ (Interpolating Cases)}

\textbf{Geometric Interpretation}:

\begin{itemize}
	\item Interpolates between $L^1$ sparsity and $L^\infty$ uniformity
	\item The unit ball has intermediate geometry
	\item Provides balance between robustness and efficiency
\end{itemize}

\textbf{Information-Theoretic Interpretation}:

\begin{itemize}
	\item Optimal for different noise models and loss functions
	\item $p = 2$: Gaussian noise (but with proper weighting)
	\item $p$ close to 1: Heavy-tailed noise, robust estimation
	\item $p$ close to $\infty$: Uniform noise, minimax estimation
\end{itemize}

\subsubsection{6.2.3 Choosing Optimal Weights}

\textbf{Fisher Information Weighting}:
For categorical distributions, the Fisher information gives natural weights:
\[
w_i = \frac{1}{p_i} \quad \text{(inverse probability weighting)}
\]

This choice makes the weighted norm scale-invariant with respect to the natural variability of each component.

\textbf{Minimum Variance Weighting}:
If we have estimates $\hat{p}_i$ with variances $\sigma_i^2$, choose:
\[
w_i = \frac{1}{\sigma_i^2} \quad \text{(inverse variance weighting)}
\]

\textbf{Importance Sampling Weighting}:
For importance sampling with proposal distribution $q$, choose:
\[
w_i = \frac{p_i}{q_i} \quad \text{(likelihood ratio weighting)}
\]

\subsection{6.3 Orlicz Norms: The Most General Suitable Framework}

\subsubsection{6.3.1 Why Orlicz Norms are Superior}

\textbf{Definition}:
\[
\|q\|_\Phi = \inf \left\{ k > 0 : \sum_{i=0}^n \Phi\left(\frac{|\text{clr}_i(q)|}{k}\right) \leq 1 \right\}
\]

\textbf{Key Advantages}:

\begin{enumerate}
	\item \textbf{Ultimate Flexibility}: Can capture any growth behavior through the Young function $\Phi$
	\item \textbf{Robust to Outliers}: Non-polynomial growth functions provide natural robustness
	\item \textbf{Information-Theoretic Optimality}: Direct connection to f-divergences and optimal estimation
\end{enumerate}

\subsubsection{6.3.2 Detailed Interpretations of Specific Orlicz Norms}

\textbf{Case 1: Exponential Orlicz Norm}
\[
\Phi(x) = e^{|x|} - 1
\]

\textbf{Geometric Interpretation}:

\begin{itemize}
	\item Extremely sensitive to large deviations
	\item Unit ball has exponentially decaying boundary
	\item Promotes extreme concentration around the origin
\end{itemize}

\textbf{Information-Theoretic Interpretation}:

\begin{itemize}
	\item Connected to exponential families and maximum entropy
	\item Optimal for sub-Gaussian distributions
	\item Related to concentration inequalities and large deviations
\end{itemize}

\textbf{Practical Applications}:

\begin{itemize}
	\item \textbf{High-frequency trading}: Extreme penalty for large position changes
	\item \textbf{Real-time systems}: Exponential penalty for delays
	\item \textbf{Safety-critical applications}: Exponential cost for failures
\end{itemize}

\textbf{Conjugate Norm}: $\Psi(y) = |y|\log|y| - |y| + 1$ (logarithmic growth)

\textbf{Case 2: Logarithmic Orlicz Norm}
\[
\Phi(x) = |x|\log(1 + |x|)
\]

\textbf{Geometric Interpretation}:

\begin{itemize}
	\item Intermediate between polynomial and exponential growth
	\item More tolerant of moderate deviations than exponential
	\item Promotes concentration but allows some spread
\end{itemize}

\textbf{Information-Theoretic Interpretation}:

\begin{itemize}
	\item Connected to entropy and mutual information
	\item Optimal for distributions with logarithmic tails
	\item Related to information-theoretic learning theory
\end{itemize}

\textbf{Practical Applications}:

\begin{itemize}
	\item \textbf{Natural language processing}: Logarithmic penalty for word frequency deviations
	\item \textbf{Network analysis}: Logarithmic cost for degree deviations
	\item \textbf{Economic modeling}: Logarithmic utility functions
\end{itemize}

\textbf{Case 3: Power-Law Orlicz Norm}
\[
\Phi(x) = \frac{|x|^p}{p} \quad \text{for } p > 1
\]

This reduces to weighted $L^p$ norms, showing that $L^p$ norms are special cases of Orlicz norms.

\textbf{Case 4: Huber-Type Orlicz Norm}
\[
\Phi(x) = \begin{cases}
\frac{x^2}{2} & \text{if } |x| \leq \delta \\
\delta|x| - \frac{\delta^2}{2} & \text{if } |x| > \delta
\end{cases}
\]
\textbf{Geometric Interpretation}:

\begin{itemize}
	\item Quadratic for small deviations, linear for large ones
	\item Provides robustness while maintaining efficiency
	\item Unit ball is a "rounded octahedron"
\end{itemize}

\textbf{Information-Theoretic Interpretation}:

\begin{itemize}
	\item Optimal for contaminated Gaussian distributions
	\item Provides robust estimation with bounded influence
	\item Connected to M-estimation and robust statistics
\end{itemize}

\textbf{Practical Applications}:

\begin{itemize}
	\item \textbf{Signal processing}: Robust to outliers and noise
	\item \textbf{Computer vision}: Robust to occlusions and artifacts
	\item \textbf{Bioinformatics}: Robust to measurement errors
\end{itemize}

\subsubsection{6.3.3 Choosing Optimal Young Functions}

\textbf{Tail Behavior Matching}:
Choose $\Phi$ to match the tail behavior of the underlying distribution:

\begin{itemize}
	\item Light tails: $\Phi(x) = e^{|x|^2}$ (super-Gaussian)
	\item Heavy tails: $\Phi(x) = |x|^p$ with $p < 2$ (sub-Gaussian)
	\item Exponential tails: $\Phi(x) = e^{|x|} - 1$
\end{itemize}

\textbf{Robustness Requirements}:

\begin{itemize}
	\item High robustness: Use $\Phi$ with linear growth for large $|x|$
	\item High efficiency: Use $\Phi$ with quadratic growth for small $|x|$
	\item Balanced: Use Huber-type functions
\end{itemize}

\textbf{Computational Considerations}:

\begin{itemize}
	\item Simple forms: $\Phi(x) = |x|^p$ for computational efficiency
	\item Smooth forms: $\Phi(x) = |x|\log(1 + |x|)$ for differentiability
	\item Separable forms: $\Phi(x) = \sum_i \phi_i(x_i)$ for decomposability
\end{itemize}

\subsection{6.4 Practical Guidelines for Norm Selection}

\subsubsection{6.4.1 Decision Framework}

\textbf{Step 1: Assess the Problem Context}

\begin{itemize}
	\item \textbf{Sparse data}: Use weighted $L^1$ norm
	\item \textbf{Robust estimation}: Use Huber-type Orlicz norm
	\item \textbf{Extreme value problems}: Use exponential Orlicz norm
	\item \textbf{Balanced performance}: Use weighted $L^2$ norm with proper weights
\end{itemize}

\textbf{Step 2: Choose Appropriate Weights}

\begin{itemize}
	\item \textbf{Equal importance}: $w_i = 1$ (unweighted)
	\item \textbf{Inverse variance}: $w_i = 1/\text{Var}(p_i)$
	\item \textbf{Fisher information}: $w_i = 1/p_i$
	\item \textbf{Domain-specific}: Choose based on application requirements
\end{itemize}

\textbf{Step 3: Validate the Choice}

\begin{itemize}
	\item \textbf{Theoretical}: Check optimality conditions
	\item \textbf{Empirical}: Cross-validation and performance testing
	\item \textbf{Robustness}: Sensitivity analysis
\end{itemize}

\subsubsection{6.4.2 Computational Algorithms}

\textbf{For Weighted $L^p$ Norms}:

\begin{itemize}
	\item \textbf{$p = 1$}: Linear programming
	\item \textbf{$p = 2$}: Quadratic programming
	\item \textbf{$p = \infty$}: Linear programming with auxiliary variables
	\item \textbf{General $p$}: Iteratively reweighted least squares
\end{itemize}

\textbf{For Orlicz Norms}:

\begin{itemize}
	\item \textbf{Gauge form}: Bisection method on $k$
	\item \textbf{Amemiya form}: Convex optimization
	\item \textbf{Proximal algorithms}: For regularization problems
\end{itemize}

\section{7. Summary}

The most suitable norms for the Aitchison geometry $(\Delta_n,\oplus;\mathbb{R},\odot)$ are:

\begin{enumerate}
	\item \textbf{Weighted $L^p$ Norms}: $\|q\|_{p,w} = \left( \sum_{i=0}^n w_i |\text{clr}_i(q)|^p \right)^{1/p}$
	\begin{itemize}
		\item Provide adaptive scaling and information-theoretic optimality
		\item Cover the range from sparse ($p=1$) to uniform ($p=\infty$) behaviors
		\item Essential for handling heteroscedastic and importance-weighted data
	\end{itemize}
	\item \textbf{Orlicz Norms}: $\|q\|_\Phi = \inf \left\{ k > 0 : \sum_{i=0}^n \Phi\left(\frac{|\text{clr}_i(q)|}{k}\right) \leq 1 \right\}$
	\begin{itemize}
		\item Provide ultimate flexibility through Young functions
		\item Enable robust estimation and tail-behavior matching
		\item Connect directly to f-divergences and information theory
	\end{itemize}
\end{enumerate}

These norms are not just mathematically valid but are the \textbf{natural} choices that respect the geometric and probabilistic structure of the probability simplex, while providing the flexibility needed for real-world applications where uniform treatment of all components is insufficient.



\end{document}